\documentclass[11pt]{article}
\usepackage[a4paper,margin=28mm]{geometry}
\usepackage{amsmath,amssymb,amsthm}
\usepackage[T1]{fontenc}
\usepackage{lmodern}
\usepackage{needspace}
\usepackage[colorlinks=true,linkcolor=blue,citecolor=blue,urlcolor=blue]{hyperref}
\newtheorem{theorem}{Theorem}
\newtheorem{lemma}{Lemma}
\newcommand{\lk}{\operatorname{lk}}
\newcommand{\rad}{\operatorname{rad}}
\newcommand{\Zact}{Z_{\mathrm{act}}}
\newcommand{\Jact}{J_{\mathrm{act}}}
\title{An elementary counterexample to a Nullstellensatz\\
for finite-type link invariants}
\author{Alec Kriebel\\
\small Independent researcher\\
\small ORCID: \href{https://orcid.org/0009-0001-9320-500X}{0009-0001-9320-500X}}
\date{12 October 2026\\\small Unrefereed research note}
\begin{document}
\maketitle
\begin{abstract}
Let \(V\) be the ordinary, uncompleted complex algebra of all finite-order
finite-type invariants of links. We give a principal ideal with a nonempty
actual-link zero set whose vanishing ideal strictly contains its radical.
The generator is \(g(4g-1)\), where \(g\) is the sum of squared pairwise
linking numbers. The obstruction is a character on the entire algebra,
obtained by evaluating the standard twist-restriction polynomial at a
nonintegral parameter. This gives a negative answer to the arbitrary-ideal
formulation of Bar-Natan's Problem~2.27, as recorded by Ohtsuki; the separate
linking-number-generated ideal question is not decided.
\end{abstract}

\section{The question and the algebra}
Bar-Natan asks whether an ideal in the algebra of finite-type link
invariants has the same radical as the ideal of invariants vanishing on its
common zeros among actual links \cite[Problem~2.27, p.~422]{Ohtsuki}.
The present note records an elementary counterexample to that literal
arbitrary-ideal assertion.

Fix \(m\geq2\), and let \(X_m\) consist of oriented, unordered
\(m\)-component links in \(S^3\), including split links. Let \(V_m\) be the
unital complex algebra of their ordinary finite-type invariants of some
finite order, with pointwise operations. An invariant has order at most
\(d\) when its positive-minus-negative crossing extension vanishes on
every link with \(d+1\) double points. Each element has a finite order;
there is no common cutoff on all elements and no infinite-order completion.
Products of invariants of orders \(d,e\) have order at most \(d+e\):
on a crossing-resolution cube their multilinear representatives have
degrees at most \(d,e\), and the product has degree at most \(d+e\).

For an ideal \(A\subset V_m\), set
\[
 \Zact(A)=\{L\in X_m:a(L)=0\text{ for all }a\in A\},\qquad
 \Jact(S)=\{v\in V_m:v|_S=0\}.
\]
Always \(\rad(A)\subseteq\Jact(\Zact(A))\), because a power
vanishing at a complex-valued evaluation forces its base to vanish there.
The same definitions apply to the algebra \(V\) on all finite-component
links, with each invariant having a finite order on that whole domain.

\Needspace{11\baselineskip}
\begin{theorem}\label{thm:main}
In \(V_m\), put
\[
 g(L)=\sum_{i<j}\lk(L_i,L_j)^2,\qquad f=g(4g-1),\qquad A=fV_m.
\]
Then \(\Zact(A)\) is the nonempty set of algebraically split links, and
\[
 g\in\Jact(\Zact(A))\setminus\rad(A).
\]
Thus \(\Jact(\Zact(A))\ne\rad(A)\). The same construction gives a
counterexample in \(V\).
\end{theorem}

\section{Polynomial restriction on the whole algebra}
The polynomial behavior of finite-type invariants on twist or braid
sequences is established prior work; see, for example,
Stanford--Trapp \cite[Theorem~E]{StanfordTrapp}, Eisermann \cite{Eisermann},
and Stoimenow \cite[Theorem~2.1, pp.~690--691]{Stoimenow}.
Stoimenow explicitly credits his earlier braiding-sequence work.
We give the particular difference argument needed here, including split
components, so that no extension of a character from a smaller algebra is
assumed.

Let \(\sigma\) be the positive generator of the two-strand braid group.
For integers \(n\geq0\), let \(L_n\) be the coherently oriented closure of
\(\sigma^{2n}\), together with \(m-2\) split unknots. Its two braid
components have linking number \(n\), all other pairwise linking numbers
vanish, and consequently \(g(L_n)=n^2\).

\begin{lemma}\label{lem:restriction}
For every \(v\in V_m\) of order at most \(d\), there is a unique
\(P_v(t)\in\mathbb C[t]\), of degree at most \(d\), satisfying
\(P_v(n)=v(L_n)\) for all \(n\geq0\).
The map \(v\mapsto P_v\) is a unital complex-algebra homomorphism.
\end{lemma}
\begin{proof}
Write \(a(n)=v(L_n)\) and \(\Delta a(n)=a(n+1)-a(n)\).
Changing one crossing in an added block \(\sigma^2\) gives
\(\sigma\sigma^{-1}=1\), by a Reidemeister~II cancellation.
Doing this in \(r\) disjoint blocks realizes
\[
 \Delta^r a(n)=\sum_{j=0}^r(-1)^{r-j}\binom rj\,v(L_{n+j})
\]
as the alternating resolution of an \(r\)-singular link. Equivalently,
the braid expression for that resolution is
\[
 \sigma^{2n}(\sigma^2-1)^r
   =\sigma^{2n+r}(\sigma-\sigma^{-1})^r.
\]
Each factor \(\sigma-\sigma^{-1}\) resolves one designated double point.
Every resolution has an even braid exponent and closes to two components;
the additional split components stay fixed.
Finite type therefore gives \(\Delta^{d+1}a(n)=0\).

The usual difference recurrence yields the Newton polynomial
\[
 P_v(t)=\sum_{r=0}^d \binom tr\,\Delta^r a(0),
 \qquad
 \binom tr=\frac{t(t-1)\cdots(t-r+1)}{r!},
\]
which agrees with \(a(n)\) at all nonnegative integers. Its uniqueness
follows because a nonzero complex polynomial cannot vanish at infinitely
many distinct integers.
The restrictions of \(u+v\), \(uv\), and \(1\) agree there with
\(P_u+P_v\), \(P_uP_v\), and \(1\), respectively. Polynomial uniqueness
proves the asserted homomorphism property.
\end{proof}

\section{The counterexample}
\begin{proof}[Proof of Theorem~\ref{thm:main}]
Pairwise linking numbers are integer-valued and have order at most one.
On any prescribed crossing-resolution cube, each is affine in the binary
resolution variables: an intercomponent crossing change changes the
relevant linking number by one up to sign, and a self-crossing change
does not change it. Its square has degree at most two on that cube.
Hence \(g\) has order at most two and \(f=4g^2-g\) has order at most four.
The sum defining \(g\) is unchanged by relabeling components or reversing
their orientations.

For every actual link, \(g(L)\) is a nonnegative integer. Thus
\[
 f(L)=0\quad\Longleftrightarrow\quad g(L)=0
 \quad\Longleftrightarrow\quad \lk(L_i,L_j)=0\text{ for all }i<j.
\]
Since \(f\) generates \(A\), this is precisely its actual-link zero set.
It contains the unlink, and \(g\) vanishes on it.

By Lemma~\ref{lem:restriction},
\[
 \chi:V_m\longrightarrow\mathbb C,\qquad \chi(v)=P_v(1/2)
\]
is a character on the entire finite-order algebra.
As \(P_g(t)=t^2\), we have
\[
 \chi(g)=\frac14,\qquad \chi(f)=0.
\]
If \(g^N=fh\) for some \(N\geq1\) and \(h\in V_m\), applying \(\chi\)
would give \(4^{-N}=0\), a contradiction.
Therefore \(g\notin\rad(A)\).

For \(V\), restrict every invariant to the same fixed-component family.
The same character proof applies; the linking-square invariant and its
order bound are defined uniformly on all finite-component links.
\end{proof}

The character is not evaluation at an actual link, since its \(g\)-value
is \(1/4\). In particular, formal evaluation at \(n=1/2\) is not
evaluation on the one-component closure of \(\sigma\).
The obstruction concerns the use of actual links as the evaluation set,
and does not contradict the classical Hilbert Nullstellensatz.

\section{Scope, attribution, and verification}
The same construction and difference proof apply to ordered, labelled
links as well. The source's separate ideal generated by the indexed
pairwise linking numbers is understood in that labelled theory.
There it differs from our principal ideal although their actual-link
zero sets coincide; its radical/vanishing-ideal question is not decided
here. The symmetric witness and the proof also apply after labels are
forgotten, as in Theorem~\ref{thm:main}. No conclusion about completed invariant algebras,
homogeneous-ideal variants, or positive characteristic is claimed.
The proof also works over any characteristic-zero field and, by pulling
back invariants to oriented links, after forgetting orientations.

The contribution recorded here is the explicit principal ideal and its
nonempty actual-link zero-set obstruction, as an elementary application of
older polynomiality. No new polynomiality method is claimed. A bounded
primary-source priority audit located no explicit equivalent earlier
counterexample in the bodies examined; it does not establish historical
firstness. The accompanying audit states its source/version access limits
and records other older routes, including the connected-sum polynomiality
of Jeong--Park \cite{JeongPark}.

The accompanying exact-rational verification checks braid resolutions,
Newton identities, linking-square symmetries, and the character calculation.
These finite controls support the written proof; they do not replace the
argument for every finite-order invariant.

\paragraph{AI and review disclosure.}
AI tools were used extensively in developing, drafting, source auditing,
and adversarially checking this result and its verification materials.
The work has not undergone conventional human peer review or refereeing
at the time of this preprint.

\begin{thebibliography}{9}
\bibitem{Ohtsuki}
T.~Ohtsuki (editor),
\emph{Problems on invariants of knots and 3-manifolds},
Geometry \& Topology Monographs \textbf{4} (volume year 2002), 377--572,
Problem~2.27, p.~422; published 1 June 2004.
\url{https://msp.org/gtm/2002/04/gtm-2002-04-024s.pdf}.
\bibitem{StanfordTrapp}
T.~Stanford and R.~Trapp,
\emph{On knot invariants which are not of finite type},
arXiv:math/9903057 (1999), Theorem~E.
\url{https://arxiv.org/abs/math/9903057}.
\bibitem{Eisermann}
M.~Eisermann,
\emph{A geometric characterization of Vassiliev invariants},
Transactions of the American Mathematical Society
\textbf{355} (2003), 4825--4846.
\bibitem{Stoimenow}
A.~Stoimenow,
\emph{Application of braiding sequences, I: On the characterization of
Vassiliev and polynomial link invariants},
Communications in Contemporary Mathematics \textbf{12}(5) (2010), 681--726.
\url{https://doi.org/10.1142/S0219199710003944}.
\bibitem{JeongPark}
M.-J.~Jeong and C.-Y.~Park,
\emph{Polynomial invariants and Vassiliev invariants},
Geometry \& Topology Monographs \textbf{4} (2002), 89--101.
\url{https://msp.org/gtm/2002/04/gtm-2002-04-007p.pdf}.
\end{thebibliography}
\end{document}
